\documentclass[conference]{IEEEtran}

\usepackage{booktabs}
\usepackage{url}
\usepackage{xspace}

\newcommand{\toolname}{Evidence-First Patch Verification\xspace}

\title{Verifiable Review of AI-Generated Patches in Real Software Projects}

\author{
\IEEEauthorblockN{Anonymous Authors}
\IEEEauthorblockA{Anonymous Institution\\
Email: anonymous@example.com}
}

\begin{document}
\maketitle

\begin{abstract}
AI coding agents increasingly produce patches for real software tasks, but a
patch that looks plausible is not necessarily safe to merge. This paper studies
patch acceptance as a verification problem: given a real task and a candidate
patch, decide whether the patch should be accepted, rejected, or escalated
based on evidence. We construct a pilot patch-verification dataset from
retained real-bug pairs, materialize source-level patch candidates, validate
labels with executable oracles, and prepare two review conditions: LLM-only
patch review and evidence-first verification. The current pre-API artifact
contains 30 validated patch candidates from 7 real-bug tasks across 2 projects,
including 9 partial-fix candidates. The pending API pilot will test whether
evidence-first verification reduces false accepts without collapsing
correct-patch recall.
\end{abstract}

\section{Introduction}

AI coding agents are increasingly used to generate patches for real codebases.
The standard question, ``does the patch pass tests?'', is insufficient when
tests are incomplete or generated changes are plausible but partial. Software
teams need a merge-gate decision: accept, reject, or escalate.

This work treats patch review as a verification problem rather than a generic
bug-finding problem. The reviewer is asked to judge a candidate patch against a
task summary and visible patch context. The evaluator uses executable or
tool-backed evidence to determine whether the decision is correct.

The motivating failure mode comes from earlier experiments in this project:
LLM-only review produced useful signals but also over-reported defects on
fixed/reference controls. That makes prompt-only review unsuitable as a merge
gate without stronger evidence discipline.

\section{Research Questions}

\textbf{RQ1.} How reliable is LLM-only review when deciding whether candidate
patches should be accepted?

\textbf{RQ2.} Can evidence-first verification reduce false accepts compared
with LLM-only review?

\textbf{RQ3.} Does evidence-first verification preserve useful correct-patch
recall, or does it only reduce false accepts by rejecting or escalating too
aggressively?

\textbf{RQ4.} Which patch types expose the most useful failure modes:
empty/no-op controls, unrelated changes, or partial fixes?

\section{Dataset Construction}

The pilot dataset is built from retained BugsInPy-derived real-bug assets. Each
source task has a buggy checkout, a fixed checkout, a task summary, touched
files, visible test hints, and retained executable oracles.

Each candidate record contains evaluator-facing fields, including patch
identity, candidate type, expected outcome, and oracle metadata. Model-visible
inputs use an anonymous candidate identifier and omit evaluator labels, oracle
paths, oracle results, and construction notes.

Patch materialization uses retained source artifacts: reference diffs between
buggy and fixed checkouts, no-op controls, applicable unrelated source changes,
and partial candidates derived from multi-hunk or multi-line reference fixes.

\section{Executable Label Validation}

Candidate labels are validated by applying each candidate patch to a copied
buggy checkout and running retained executable oracles. This guards against
purely prompt-based or manually asserted labels. Correct patches are expected
to pass all retained oracles. Negative candidates are expected to fail at least
one retained oracle.

\section{Review Conditions}

\subsection{LLM-Only Review}

The reviewer sees the task summary, visible context, and candidate patch. It
must output one JSON object with decision, confidence, claims, rationale, and
uncertainty. It does not see hidden oracle paths or oracle results.

\subsection{Evidence-First Verification}

The reviewer sees the same task and patch context plus model-visible evidence
source metadata and visible test hints. It must tie every accept/reject
decision to concrete visible evidence. If the visible evidence is insufficient,
it should escalate.

\subsection{Oracle Upper Bound}

The evaluator can use hidden labels and oracle outcomes to produce an upper
bound. This is not model capability and must be reported separately.

\section{Metrics}

Primary metrics are patch-level: false accept rate, false reject rate, accepted
precision, correct-patch recall, escalation rate, invalid-output rate, and cost.
Escalation is neither accept nor reject. It should be reported separately
because reducing false accepts by escalating everything is not a useful merge
gate.

\section{Current Pre-API Results}

The current no-API baselines validate the metric implementation and expected
tradeoffs. These baselines are not model-review results. They expose the
merge-gate tension: accepting everything preserves correct-patch recall while
accepting all wrong patches, and rejecting everything avoids false accepts
while rejecting all correct patches.

% Generated pre-API paper tables. Requires booktabs.

\begin{table}[t]
\centering
\caption{Dataset by project.}
\label{tab:dataset-projects}
\begin{tabular}{lr}
\toprule
Project & Count \\
\midrule
httpie & 22 \\
luigi & 8 \\
\bottomrule
\end{tabular}
\end{table}


\begin{table}[t]
\centering
\caption{Candidate types.}
\label{tab:candidate-types}
\begin{tabular}{lr}
\toprule
Candidate type & Count \\
\midrule
buggy\_noop & 7 \\
correct\_reference & 7 \\
irrelevant\_patch & 7 \\
partial\_fix & 9 \\
\bottomrule
\end{tabular}
\end{table}


\begin{table}[t]
\centering
\caption{Executable validation summary.}
\label{tab:validation-summary}
\begin{tabular}{lr}
\toprule
Item & Value \\
\midrule
Records & 30 \\
Patch applied & 30 \\
Oracle ran & 30 \\
Oracle all passed & 7 \\
All validated & true \\
\bottomrule
\end{tabular}
\end{table}


\begin{table*}[t]
\centering
\caption{No-API baseline metrics. These baselines validate metric behavior but are not model-review results.}
\label{tab:no-api-baselines}
\begin{tabular}{lrrrrrr}
\toprule
Baseline & Accepted precision & False accept & Correct recall & False reject & Escalation & Invalid output \\
\midrule
accept-all & 0.2333 & 1.0000 & 1.0000 & 0.0000 & 0.0000 & 0.0000 \\
oracle upper bound & 1.0000 & 0.0000 & 1.0000 & 0.0000 & 0.0000 & 0.0000 \\
reject-all & -- & 0.0000 & 0.0000 & 1.0000 & 0.0000 & 0.0000 \\
\bottomrule
\end{tabular}
\end{table*}


\begin{table}[t]
\centering
\caption{Deterministic no-API reproducibility check.}
\label{tab:reproducibility}
\begin{tabular}{lr}
\toprule
Item & Value \\
\midrule
Matched & true \\
Checked files & 7 \\
Mismatches & 0 \\
Missing & 0 \\
\bottomrule
\end{tabular}
\end{table}


\section{API Pilot Plan}

The first API pilot should run only two conditions on the validated candidates:
LLM-only review and evidence-first verification. The same model must be used
for both conditions. Before API calls, the executor must create ignored local
configuration files, document the concrete OpenRouter model slug and selection
rationale, run strict preflight, and complete a two-candidate smoke run before
the full run.

\textbf{Pending results.} This draft intentionally leaves the real API result
section blank. It must not be filled until real \texttt{reviews.jsonl},
\texttt{metrics.json}, failure examples, and a stop/continue gate report are
available.

\section{Reproducibility and Handoff Controls}

The pre-API artifact includes deterministic reproduction checks for local
dataset construction. The original no-API pilot and a reproduced run are
compared by hashing deterministic output files. Runtime work directories, raw
API responses, external checkouts, and environment-dependent files are not
treated as deterministic reproducibility evidence.

The execution plan also separates local pre-API checks from model experiments.
The current handoff packet refreshes readiness, paper readiness, plan progress,
human-required inputs, Git-sync state, OpenRouter catalog visibility, and
deterministic reproducibility.

\section{Model Selection Boundary}

The first real API pilot is a within-model comparison. The current shortlist
recommends a concrete OpenRouter-accessible model candidate, but this is not an
experiment decision by itself. The user must confirm the slug and rationale
before local configuration is generated. The first pilot controls for base
model capability but cannot support cross-model generality.

\section{Threats to Validity}

Dataset size is small. The current pilot is designed to validate the method and
failure surfaces, not to make broad claims about all AI-generated patches. The
partial-fix candidates are source-backed and oracle-checkable, but they are
constructed from retained reference diffs rather than generated by live coding
agents. A later stage should add model-generated patches or SWE-bench-style
tasks.

Visible test hints may not represent the evidence actually available in all
engineering workflows. Model behavior can drift over time and across providers.
Every API run must record model slug, provider, date, prompt version, decoding
settings, cost, raw response path, and invalid-output status.

\section{Conclusion}

The current artifact establishes a validated patch-verification pilot and a
ready-to-run API protocol. It does not yet establish that evidence-first
verification improves over LLM-only review. That claim depends on the pending
API pilot and must be evaluated using false accept rate, accepted precision,
correct-patch recall, escalation rate, and invalid-output rate.

\end{document}
